\section{Subgraph Retrieval and Traversal}
\label{sec:gr-subgraph}

In Graph RAG the retrieval target is not a list of independent passages but a
relevant \emph{subgraph} of connected entities and relations, which is then
handed to the generator as structured context~\cite{peng2024graphrag}. A common
strategy is to first identify seed nodes---for instance the entities mentioned in
the query, located by dense similarity---and then traverse outward along edges to
gather the neighbourhood of facts that surround them~\cite{peng2024graphrag}.
LightRAG operationalizes this with a dual-level retrieval scheme that combines
low-level retrieval of specific entities with high-level retrieval of broader
relational context, so that both fine-grained and thematic information are
captured in a single pass~\cite{guo2024lightrag}. Retrieving a connected subgraph
rather than isolated chunks is precisely what lets the model see how the retrieved
facts relate to one another~\cite{peng2024graphrag}.
